\documentclass[10pt,a4paper]{amsart}
\usepackage[margin=19mm]{geometry}
\usepackage{amsmath,amssymb,mathtools}
\usepackage[scr=rsfs,cal=euler]{mathalfa}
\usepackage{xfrac}
\usepackage[unicode,hidelinks,bookmarks=false]{hyperref}
\newcommand{\ie}{i.e.\ }
\newcommand{\cf}{cf.\ }
\newcommand{\resp}{respectively\ }
\newcommand{\Subsection}{\subsection}
\providecommand{\nobreakdash}{\nobreak\hbox{-}}
\setlength{\emergencystretch}{2em}
\multlinegap0pt
\usepackage{bm}
\usepackage{bbm}
\usepackage{dsfont}
\usepackage{xspace}
\usepackage{cancel}
\usepackage{mathtools}
\usepackage{amsthm}
\makeatletter
\@ifundefined{theorem}{\newtheorem{theorem}{Theorem}}{}
\makeatother
\DeclareMathOperator{\Z}{\mathbb{Z}}
\title[Global Hypoellipticity and Solvability with Loss of Derivati]{Global Hypoellipticity and Solvability with Loss of Derivatives on the Torus}
\author{Andr\'e Pedroso Kowacs, Alexandre Kirilov}
\date{Source: arXiv:2503.17466v3 (CC BY 4.0)}
\begin{document}
\maketitle

\noindent\emph{Source and licence.} Global Hypoellipticity and Solvability with Loss of Derivatives on the Torus. Version arXiv:2503.17466v3 (CC BY 4.0), distributed under the Creative Commons Attribution 4.0 International licence, as recorded in the arXiv licence metadata for this exact version and shown on the abstract page https://arxiv.org/abs/2503.17466v3. Full rights notice: licenses/arxiv-2503.17466-CC-BY-4.0.txt.\par\smallskip

\noindent\textit{Editorial scope.} Counted unit: the Theorem whose source label is \texttt{teoGS-closed} in the article (source0.tex lines 928--934, source label \texttt{teoGS-closed}), delivered together with the article's own complete original proof of it (source0.tex lines 936--1025, one \texttt{proof} environment of 90 lines). No mathematics is written, repaired or re-derived: statement and proof are the article's own text line for line. Required context retained uncounted: theoGS (lines 708--714). Every definition, display and label of those lines is retained unchanged. Omitted from this package and not redistributed: 1--707, 715--927, 935--935, 1026--1708. Nothing inside a retained excerpt is omitted or altered: every retained body is delivered exactly as the article's lines are written, so no omission receipt of any kind accompanies this package. Citations: the retained excerpts carry no citation command at all, so no bibliography entry of the article is transcribed and no reference is redistributed. Reference adaptation: none was needed and none was made; every reference inside the delivered unit resolves inside it, so no unresolved reference or citation is produced. Printed slips kept literal: no printed word, symbol, hypothesis or number of the counted statement or of its proof was altered. Layout and class: the article's own class is \texttt{amsart} with packages including inputenc, amsmath, amssymb, amsthm, dsfont, mathrsfs, xcolor, hyperref, mathtools; the wrapper is the lane's pinned \texttt{amsart} editorial wrapper at 10pt on A4, which restates the theorem-like environments the retained text needs and adds the article's own macro definitions that the retained text needs (\texttt{\textbackslash Z}), with extra wrapper packages (bm, bbm, dsfont, xspace, cancel, mathtools). The delivered numbering is the unit's own. Assets: no retained span contains a figure environment, a graphics-inclusion command, a PDF-page inclusion, a table or a TikZ picture; no image file is copied into this package, the package declares no asset dependency and no image file is redistributed with it. Licence: version arXiv:2503.17466v3 of this article is distributed under the Creative Commons Attribution 4.0 International licence, as recorded in the arXiv licence metadata for this exact version and shown on the abstract page https://arxiv.org/abs/2503.17466v3. A full rights notice is delivered at licenses/arxiv-2503.17466-CC-BY-4.0.txt. Characters: every retained excerpt is pure ASCII, so the delivered engine, XeLaTeX with its default Latin Modern text font, reports no missing character.
	\begin{theorem}\label{theoGS}
		Let \( r \geq 0 \) and \( p \in S^{[m]}(\mathbb{Z}^n) \). The operator \( p(D) \) is GS-$r$ if and only if there exists a constant \( K > 0 \) such that
		\begin{equation}\label{ineq-GS}
			|p(\xi)|\geq K|\xi|^{m-r},
		\end{equation}
		for every \( \xi \in \mathbb{Z}^n \setminus \{0\} \) such that \( p(\xi) \neq 0 \).
	\end{theorem}

	\begin{theorem}\label{teoGS-closed}
		Let \( p \in S^{[m]}(\mathbb{Z}^n) \) and \( r \geq 0 \). The operator \( p(D) \) is GS-\( r \) if and only if for every \( k \in \mathbb{R} \), the mapping
		\[
		p(D): H^{k + m - r}(\mathbb{T}^n) \to H^k(\mathbb{T}^n)
		\]
		has closed range.
	\end{theorem}

		\begin{proof}
		Suppose that $p(D)$ is GS-$r$. Then by Theorem \ref{theoGS} we have that there exists $K>0$ such that
		\begin{align}\label{ineq-GS-strong}
			|p(\xi)|\geq K|\xi|^{m-r},
		\end{align}
		for every $\xi\in\Z^n\backslash\{0\}$ such that $p(\xi)\neq 0$.
		
		First notice that for $m'>m$, $p(D)$ if of order $m'$. Therefore $p(D)$ maps $H^{k + m - r}(\mathbb{T}^n)\subset H^{k+m'}(\mathbb{T}^n)$ to $H^k(\mathbb{T}^n)$, and the mapping \(p(D): H^{k + m - r}(\mathbb{T}^n) \to H^k(\mathbb{T}^n)\) is well defined. Now Let $ u \in H^{k + m - r}(\mathbb{T}^n) $. Then 
		\begin{align*}
			\|p(D)u\|_{H^k(\mathbb{T}^n)}^2 &= \sum_{\xi \in \mathbb{Z}^n} (1 + \|\xi\|^2)^k \left| \widehat{p(D)u}(\xi) \right|^2 \\
			&= \sum_{\xi \in \mathbb{Z}^n} (1 + \|\xi\|^2)^k |p(\xi)|^2 |\widehat{u}(\xi)|^2 \\
			&\geq K(1+n)^{|m-r|} \sum_{\xi \in \mathbb{Z}^n} (1 + \|\xi\|^2)^{k + m - r} |\widehat{u}(\xi)|^2 \\
			&= K(1+n)^{|m-r|} \|u\|_{H^{k + m - r}(\mathbb{T}^n)}^2,
		\end{align*}
		
		Since both $ H^k(\mathbb{T}^n) $ and $ H^{k + m - r}(\mathbb{T}^n) $ are Banach spaces, it follows that $ p(D):H^{k + m - r}(\mathbb{T}^n)\to H^k(\mathbb{T}^n)$ has closed range. 
		
		
	Conversely, suppose that $ p(D) $ is not GS-$ r $. Then, by Theorem \ref{theoGS}, inequality \eqref{ineq-GS-strong} does not hold for any $ K > 0 $. Consequently, there exists a sequence of distinct elements $ \xi_j \in \mathbb{Z}^n \setminus \{0\} $ such that
	\begin{equation}\label{ineq_contradiction_closed}
		0 < |p(\xi_j)| \leq \frac{1}{j} |\xi_j|^{m - r}, \quad  j \in \mathbb{N}.
	\end{equation}
	
	Define a sequence of distributions $ (u_\ell)_{\ell \in \mathbb{N}} \in \mathscr{D}'(\mathbb{T}^n) $ whose Fourier coefficients are given by:
	\[
	\widehat{u_\ell}(\xi) =
	\begin{cases}
		(1 + \|\xi_j\|^2)^{(-k - m + r)/2}, & \text{if } \xi = \xi_j, \, j \leq \ell, \\
		0, & \text{otherwise.}
	\end{cases}
	\]
	
	For each $ \ell \in \mathbb{N} $, the Fourier coefficients of $ p(D)u_\ell $ are:
	\[
	\widehat{p(D)u_\ell}(\xi) =
	\begin{cases}
		p(\xi_j)(1 + \|\xi_j\|^2)^{(-k - m + r)/2}, & \text{if } \xi = \xi_j, \, j \leq \ell, \\
		0, & \text{otherwise.}
	\end{cases}
	\]
	
	Since both $ p(D)u_\ell $ and $ u_\ell $ have only finitely many nonzero Fourier coefficients, it follows that $ p(D)u_\ell, u_\ell \in \bigcap_{s \in \mathbb{R}} H^s(\mathbb{T}^n) = C^\infty(\mathbb{T}^n) $.
	
	Now consider the sequence of Fourier coefficients defined by:
	\[
	\widehat{f}(\xi) =
	\begin{cases}
		p(\xi_j)(1 + \|\xi_j\|^2)^{(-k - m + r)/2}, & \text{if } \xi = \xi_j, \, j \in \mathbb{N}, \\
		0, & \text{otherwise.}
	\end{cases}
	\]
	
	We claim that these coefficients define a distribution $ f \in H^k(\mathbb{T}^n) $. Indeed, we compute:
	\begin{align*}
		\|f\|_{H^k(\mathbb{T}^n)}^2 &= \sum_{\xi \in \mathbb{Z}^n} (1 + \|\xi\|^2)^k |\widehat{f}(\xi)|^2 \\
		&= \sum_{j \in \mathbb{N}} (1 + \|\xi_j\|^2)^k |p(\xi_j)|^2 (1 + \|\xi_j\|^2)^{-k - m + r} \\
		&\leq (1 + n)^{|m - r|} \sum_{j \in \mathbb{N}} \frac{1}{j^2} < \infty.
	\end{align*}
	
	Thus, $ f \in H^k(\mathbb{T}^n) $.
	
	\smallskip
	Next, observe that $ p(D)u_\ell \to f $ in $ H^k(\mathbb{T}^n) $ as $ \ell \to \infty $. Indeed, we have:
	\[
	\|p(D)u_\ell - f\|_{H^k(\mathbb{T}^n)}^2 = \sum_{j \geq \ell} (1 + \|\xi_j\|^2)^k |p(\xi_j)|^2.
	\]
	
	Since $ \sum_{j \geq 1} (1 + \|\xi_j\|^2)^k |p(\xi_j)|^2 $ converges, it follows from the Cauchy criterion for series that:
	\[
	\sum_{j \geq \ell} (1 + \|\xi_j\|^2)^k |p(\xi_j)|^2 \to 0 \quad \text{as } \ell \to \infty.
	\]
	
	Hence, $ \|p(D)u_\ell - f\|_{H^k(\mathbb{T}^n)}^2 \to 0 $, which implies $ p(D)u_\ell \to f $ in $ H^k(\mathbb{T}^n) $.
	
	Finally, note that $ f $ is not in the range of $ p(D): H^{k + m - r}(\mathbb{T}^n) \to H^k(\mathbb{T}^n) $. Indeed, if $ p(D)u = f $, then comparing Fourier coefficients yields:
	\[
	\widehat{u}(\xi) =
	\begin{cases}
		(1 + \|\xi_j\|^2)^{(-k - m + r)/2}, & \text{if } \xi = \xi_j, \, j \in \mathbb{N}, \\
		0, & \text{otherwise.}
	\end{cases}
	\]
	
	However, this implies that $ u \notin H^{k + m - r}(\mathbb{T}^n) $, since:
	\[
	\|u\|_{H^{k + m - r}(\mathbb{T}^n)}^2 = \sum_{j=1}^\infty (1 + \|\xi_j\|^2)^{(k + m - r)} (1 + \|\xi_j\|^2)^{(-k - m + r)} = \sum_{j=1}^\infty 1 = \infty.
	\]
	
	Therefore $ p(D): H^{k + m - r}(\mathbb{T}^n) \to H^k(\mathbb{T}^n) $ does not have closed range.
  \end{proof}

\end{document}
